RQ1 asks whether the selector captures the known visible effects in the tested
mutation classes and independent reference sample. RQ2 measures how much native
capture work and elapsed time it avoids. RQ3 examines the reproducibility and
investigability of retained findings, including failures and rendering variation.

\subsection{Corpus identity and outcome accounting}
The evaluation treats Data.gov, Destatis and Enron as distinct populations.
Archive checksums, member names and duplicate occurrences identify each selected
corpus version. Legacy XLS conversion is reported separately from XLSX cleanup.
The ongoing Enron fallback selection contains 20,872 inputs, of which 20,503 have
prepared artifacts and 369 have unsuccessful preparation outcomes. Its visual
scope contains 16,799 workbooks. These are preparation counts for the current
selection, not final visual-validation results.

A dated evaluation cutoff will report compared workbooks, differences, failures,
unsupported representations and unfinished work. Repeated attempts have their own
counts. An unchanged package can bypass visual work under an exact input-identity
check, but a failed capture cannot contribute an equality observation. The full
recursive Enron XML-derived corpus is a separately identified population.

\subsection{Sampling and reference controls}
The random evaluation sample is selected from a frozen eligible frame with a
recorded seed before inspecting visual findings. Strata use corpus identity,
package size and structural features available before validation. Low-reduction
and unchanged files remain represented. A separately labelled development stress
set contains large and historically problematic workbooks.

Feasible sampled files receive full native reference captures over the relevant
visual area using a capture plan independent of the candidate selector. A
full-reference timeout remains in the denominator. It is not replaced silently
by an easier workbook. Mutations with known visible effects cover cell values,
styles, geometry, conditional dependencies and objects. Mutations that Excel
normalizes away are reported separately.

\subsection{Detection and cost metrics}
A selector miss is a native reference difference outside the selected coverage.
Report detection by mutation class and by completed reference workbook. A changed
capture is counted separately from a confirmed transformation regression.
Unsupported references, failed geometry reads and missing partner images retain
explicit coverage outcomes.

Cost measurements separate XML planning, geometry, native rendering, transfer,
image comparison and persistence where instrumentation permits. Report image
count, compared pixels and retained bytes alongside elapsed time. Batch and
concurrency experiments use matched inputs and repeated runs. A changing
production queue does not identify the causal benefit of a tuning parameter.
Full-corpus cost estimates use measured visible geometry and stratified samples,
with uncertainty and assumptions about hidden regions and object extents.

\subsection{Current verification boundary}
The portable-viewer smoke experiment retained a deliberately different pair,
loaded its images through relative static assets and displayed its diagnostics
without a running control API. An empty-partition test verified explicit selected
and excluded counts without opening Excel. These tests verify execution and
reporting boundaries. They do not establish selector recall.

\textit{Evaluation freeze pending.} Complete the independent mutation matrix,
seeded native reference sample and matched cost tables before making detection
or speedup claims. Retain unresolved development cases and the original pixel
observations when a later retry or software fix changes their classification.
